problem:  
Let \( M^4 \) be a compact, simply connected Riemannian 4–manifold with nonnegative sectional curvature and an effective isometric circle action \( S^1 \curvearrowright M \). The orbit space \( M^* = M/S^1 \) is a 3–dimensional Alexandrov space with nonnegative curvature, whose boundary corresponds to the fixed-point set of the action. In all known examples (e.g. \(M \cong S^4, \mathbb{C}P^2, S^2 \times S^2, \mathbb{C}P^2 \# \pm \mathbb{C}P^2\)), \( M^* \) is homeomorphic to a 3–sphere or a 3–ball, and the image of the singular orbits—denoted by \(\Sigma^* \subset M^*\)—forms an **unknotted** embedded 1–dimensional graph (an arc or a circle).  

**Problem:**  
Assume that the orbit space \(M^*\) of such a manifold is homeomorphic to \(S^3\), and that \(\Sigma^*\subset S^3\) is the image of the singular orbits. Is it possible for \(\Sigma^*\) to be a **knotted** embedded circle in \(S^3\)?  
Equivalently:
> Does there exist a compact, simply connected, nonnegatively curved 4–manifold \( (M^4, g) \) with an isometric circle action whose orbit space is \(S^3\) and whose singular locus projects to a nontrivial knot in \(S^3\)?  

Alternatively, prove that the nonnegative curvature condition forces \(\Sigma^*\) to be unknotted.

Why is it a "good" problem:  
This problem makes a precise inquiry into the interplay between curvature, symmetry, and topological knotting in 4–manifolds. It refines the existing classification program (Hsiang–Kleiner; Fintushel; Grove–Searle; Grove–Wilking) by asking whether metric curvature conditions impose **topological unknotting** constraints on the singular orbit structure. The phenomenon in question—the potential existence of knotted singular sets—has never been observed, yet no general obstruction is known.  

From a conceptual standpoint:
- It tightly couples **3‑dimensional Alexandrov geometry** (orbit space of nonnegative curvature) with **knot theory** and **4‑dimensional transformation groups**; resolving it could reveal curvature-based topological rigidity phenomena.  
- If an example exists, it would expand the currently known roster of nonnegatively curved 4‑manifolds; if impossible, it would uncover a new bridge between geometric curvature bounds and knot-theoretic simplicity.  
- The formulation is elementary (is there a knotted singular orbit?) yet its resolution requires deep geometric insight joining Riemannian, Alexandrov, and topological techniques.  

Hence the problem fulfills the “narrow entrance, profound depth” criterion: its statement is simple and tangible, but its resolution could reshape our understanding of symmetry and curvature in dimension four.